(1) 轴上焦点到距光轴 $r$ 的波带片点的距离为 $\sqrt{f^2+r^2}$，相对轴中心路程差为 $\Delta\ell=\sqrt{f^2+r^2}-f\simeq r^2/(2f)$，要求 $r\ll f$。将第 $m$ 个半波带边界设为 $\Delta\ell=m\lambda/2$，得到 $r_m^2\simeq m\lambda f$。
(2) 相邻半波带的光程差增加 $\lambda/2$，在焦点处两相邻带的复振幅相差 $\pi$，天然趋于抵消。若把偶数半波带遮挡、保留奇数带，留下的邻近透光带间相位相差 $2\pi$，因而在指定焦点方向可以相干增强。
(3) 令 $m=1$ 可得 $r_1^2=\lambda f$，故焦距 $f=r_1^2/\lambda$；其余带界为 $r_m=\sqrt m\,r_1$，与 $f$ 的几何关系一致。结果显示波带宽度随半径增大逐渐变窄，制造时必须保持边界位置精确。
(4) 波带片图案固定时 $r_1$ 不变，故 $f(\lambda)\propto1/\lambda$，波长越长焦距越短，与常见普通玻璃折射透镜的正常色散趋势不同。若波带不再满足 $r/f\ll1$，应改用精确方程 $\sqrt{f^2+r_m^2}-f=m\lambda/2$。
